% ID: SEL-D0080-OA-E01
% ESTADO: cotejada. Segunda pasada visual y borrador compilado revisados.
% Fuente DOC: Selectividad/sel_2012_matematicas_aplicadas/2 MATEMÁT. APLIC. CC. SS/EXAMEN 5-ANDALUCÍA.doc
% SHA-256 original: 87b729cdd4632386690a9390bdfa17dd7d5c902a2cf7e5801feea45e2b181969
% Página convertida: 1; opción A; ejercicio 1.

\textbf{(2.5 puntos)} Un comerciante dispone de 1200 euros para comprar dos tipos de manzanas A y B. Las del tipo A las compra a $0.60$ euros/kg y las vende a $0.90$ euros/kg, mientras que las del tipo B las compra a 1 euro/kg y las vende a $1.35$ euros/kg.

Sabiendo que su vehículo a lo sumo puede transportar 1500 kg de manzanas, \textquestiondown{}cuántos kilogramos de cada tipo deberá adquirir para que el beneficio que obtenga sea máximo? \textquestiondown{}Cuál sería ese beneficio?
